Considera els punts $A(0,1,-2)$ i $B(2,-3,4)$.

\begin{apartats}

\apartat[1,25]{0,75}
Escriu les equacions vectorial i paramètriques de la recta $r$ que passa per $A$ i $B$.

\begin{solucio}
$\overrightarrow{AB}=(2,-4,6)$, i també serveix $\vec v=(1,-2,3)$, proporcional.\\
Vectorial: $(x,y,z)=(0,1,-2)+t(1,-2,3)$.\quad
Paramètriques: $\left\{\begin{aligned}x&=t\\y&=1-2t\\z&=-2+3t\end{aligned}\right.$
\end{solucio}

\begin{tria}{punts-recta}
\itemtria{alineat}{1}{1,25}
Troba el valor de $a$ perquè el punt $C(a,-7,10)$ estigui alineat amb $A$ i $B$.

\begin{solucio}
$C$ ha de ser de la recta $r$. De la segona coordenada, $1-2t=-7$, $t=4$. Amb $t=4$, $z=-2+12=10$, que
coincideix: el punt és de la recta si $a=x=4$.
\end{solucio}

\itemtria{talls-plans-coordenats}{1}{1,25}
Troba els punts on la recta $r$ talla els plans coordenats $z=0$ i $x=0$.

\begin{solucio}
Amb $z=0$: $-2+3t=0$, $t=\frac23$, i el punt és $\left(\frac23,-\frac13,0\right)$.\\
Amb $x=0$: $t=0$, i el punt és $A(0,1,-2)$ mateix: la recta talla el pla $x=0$ a $A$.
\end{solucio}
\end{tria}

\begin{nomesllarg}
\apartat{0,75}
Escriu la recta $r$ en forma contínua i en forma implícita.

\begin{solucio}
Contínua: $\dfrac{x}{1}=\dfrac{y-1}{-2}=\dfrac{z+2}{3}$.\\
De $-2x=y-1$ i $3x=z+2$:
$\left\{\begin{aligned}2x+y-1&=0\\3x-z-2&=0\end{aligned}\right.$
\end{solucio}
\end{nomesllarg}

\end{apartats}
